% Transcription — fonds Alexandre Grothendieck, shelfmark 6, batch 1 (pages 1-20).
% Established from the facsimile archives/batches/6/batch-01.pdf, cut from a
% copy of 6.pdf fetched through the project's own relay
% (https://grothendieck-relay.onrender.com/source/6.pdf) rather than from
% grothendieck.umontpellier.fr directly; per the skill transcribe-grothendieck.
% The batch file carries no cover sheet: PDF page k is the archivists' page k,
% and their pencil numbers bottom-left agree throughout. The folder is
% sixty-six pages; this is the first of four batches (the others are
% transcribed in parallel by other passes).
%
% Pass: Opus 5.5 (claude-opus-5-5), 2026-09-26 — first pass, unchecked against the pages by a human.
%
% The folder sits in the inventory group "4-9" « Cristaux »; it has its own
% date in the catalogue, which \dating{} carries verbatim, with no group
% sentence.
%
% Pages skipped: 2, the blank verso of the first leaf (only the printed
% letterhead of the Istituto Matematico of Pisa and show-through; nothing
% of his).
%
% Pages summarised: none.
%
% The letter is HIS OWN and is transcribed in full. Evidence: pages 1 and
% 3-20 are a typescript on the letterhead paper of the Università di Pisa,
% Istituto Matematico « Leonida Tonelli », opening « Cher John, » and written
% throughout in the first person (« je n'ai pas écrit de démonstration »,
% « mon papier bleu sur De Rham », « Etant loin de mes bases »); page 1 is
% dated at the top right in his hand, in ink, « Pise, Mai 1966 »; and the
% typescript carries dozens of corrections and insertions in the same ink
% hand (e.g. « et passage au dual. » p. 10, « Effectivement, ça ne marche
% pas ! » in the margin of p. 13, the boxed addition on Hom ⊗ Q at the foot
% of p. 8, « (et "est" par "définit") » p. 18), plus pencil marginalia
% (« marche en car. 0, et pas en car p>0 », p. 17) — i.e. this is his own
% typed copy of the May 1966 letter to Tate, corrected by him; the typed
% text and his ink and pencil interventions are given in full. The letter's own
% pagination (top right, from p. 1 to p. 19 in this batch) runs one behind
% the archivists' numbering, from page 3 on; it is recorded in a \note at
% each page. The letter continues into batch 2 (it breaks off mid-word at
% « condi- », §2.3).
%
% Apparatus: ink insertions (and typed interlinear insertions called in by an
% ink stroke) are \add{}; words struck at the machine or in ink are
% \struck{} when legible, \struck{\ill{}} when not; machine overstrikes of a
% single mistyped word are silently read through. Underlined words are
% \emph{}; letters underlined at the machine (Z, F, M, O, R, C, Ω) are kept as
% \underline{} in math. Obvious slips of the typescript are kept with a
% \note{sic}: « De Rahm » (p. 1), « isogristaux » (p. 9), « voisinage
% infinitésimal de R dans S » (p. 4), F_R for F_S (p. 5), poids 1 for 2
% (p. 8), the missing p in VF = id (p. 8), C_{S^an} (p. 19).
%
% What the batch is about. Chapter 1, « La notion de cristal » (1.0-1.17):
% the category C of nilpotent thickenings of S over R; crystals of modules
% and, generally, F-crystals as cartesian sections (1.1); functoriality and
% inverse images (1.2-1.3); the value on S, not faithful in general (1.4);
% examples: perfect k and Witt/Cohen rings (1.5), S étale over R (1.6), S a
% subscheme of R and formal completion (1.7); n-connexions,
% pseudo-stratifications and stratifications, with the equivalence crystals
% = stratified objects for S formally smooth over R (1.8); characteristic
% 0 (curvature) versus Cartier in char. p (1.9); p-crystals via frobenius,
% p^{-1}-crystals, bicrystals (« cristaux de Dieudonné ») with FV = p,
% bicrystals of weight i and the Tate bicrystal (1.10); isocrystals,
% p-isocrystals and bi-isocrystals, the fully faithful Isbicr(S,i) ->
% Isbicr(S) (1.11); the perfect-field case, Dieudonné modules and
% effective bi-isocrystals T^{-i} ⊗ N (1.12); S smooth over perfect k:
% lifting to W_n, stratified modules on a lifting X, the transcendental
% approach via the fundamental group of X_C, rigid-analytic X_K and
% noetherianness (1.13); explicit frobenius via a lifting f_X (1.14);
% bicrystals over an arbitrary base (1.15); the correction of 1.0 to
% open U ⊂ S and the « site cristallogène » with its topos Topcr_{S/R}
% (1.16); the cohomological test that the crystalline cohomology of O_C
% should give De Rham (1.17). Chapter 2, « La cohomologie de De Rham est un
% cristal »: evidence from Washnitzer-Monsky and Gauss-Manin, the spectral
% sequence of his « papier bleu » (2.1); the desired complex of crystals
% DR(f) with base change, Künneth, duality into T^d[2d] and DR(f)(S) ≃
% Rf_*(Ω•_{X/S}) (2.2); the frobenius F_p on DR(f)_p (2.3, continued in
% batch 2).
%
% Notation fixed once for the batch: machine-underlined letters as
% \underline{Z}, \underline{F}, \underline{M}, \underline{O}_T,
% \underline{\Omega}; frob_S, M^{(p)}, f_W, W_n, S_n, X_K; Isbicr, IsBicr,
% Bicr, Topcr as roman operator names; « Specf » as \operatorname{Specf};
% his ink calligraphic R as \mathcal{R} (p. 13); the arrow ~> of the
% typescript (a functor on objects) as \rightsquigarrow.
%
% Readings to check first: « défini » p. 5 (ink over the typed word); the
% prefix of « bi-isocristaux » and the struck ink word before « pleinement
% fidèle » p. 9; « ne marche » in the margin of p. 13; « sur » and
% « ouverte » in the ink addition at the foot of p. 16; « inverse » in
% « au sens inverse » p. 17; the pencil marginal « (faut d'introduire des
% puiss. div.) » p. 17, read largely by conjecture.

\documentclass[11pt,a4paper]{article}
% Le préambule commun à toutes les transcriptions du fonds Grothendieck.
%
% Il tient en une idée : **l'incertitude se compose, elle ne se résout pas.**
% Une transcription qui lisse un mot douteux a détruit la seule chose pour
% laquelle on la faisait — un lecteur doit pouvoir savoir, à chaque endroit,
% ce qui était sur le feuillet et ce qui a été ajouté. Les six macros qui
% suivent sont donc l'appareil critique, et non de la décoration.
%
% Les mêmes commandes sont reconnues par `scripts/render.mjs`, qui produit la
% vue de lecture du site. Une macro ajoutée ici doit l'être aussi là-bas,
% faute de quoi le rendu échouera bruyamment — ce qui est voulu : un rendu
% silencieusement faux est pire qu'un rendu absent.

\ProvidesPackage{grothendieck}[2026/08/08 Transcriptions du fonds Grothendieck]

\usepackage{amsmath,amssymb,amsthm}
\usepackage{mathrsfs}
\usepackage{stmaryrd}
% Les diagrammes commutatifs sont partout dans ce fonds : c'est la langue
% même de la géométrie algébrique de Grothendieck, et une transcription qui
% les rendrait en prose ne transcrirait rien.
\usepackage{tikz-cd}
% Français par défaut — c'est la langue des feuillets — mais l'anglais est
% chargé aussi : la traduction et le résumé sont des documents anglais, et la
% typographie française (espaces avant les deux-points) y serait une faute.
% Ces documents basculent par \selectlanguage{english} en tête de corps.
\usepackage[english,french]{babel}
\usepackage{fontspec}
\usepackage{geometry}
\geometry{a4paper,margin=25mm,marginparwidth=28mm}
\usepackage{marginnote}
\usepackage{xcolor}
% ulem, not soul. `soul`'s \st and \ul are built on a character-by-character
% scan that XeLaTeX's font machinery breaks: the document fails with an
% undefined control sequence rather than degrading. ulem does the same two jobs
% and is Unicode-engine safe. `normalem` stops it hijacking \emph, which the
% transcriptions use for Grothendieck's own emphasis.
\usepackage[normalem]{ulem}
\usepackage{hyperref}
\hypersetup{colorlinks=true,linkcolor=black,urlcolor=black,pdfborder={0 0 0}}

% --- Métadonnées du lot -----------------------------------------------------
% Reprises telles quelles par le script de rendu, qui les affiche en tête de
% la vue de lecture. Elles fixent ce que le fichier prétend couvrir : sans
% elles, une transcription flotte sans point d'attache dans un fonds de seize
% mille feuillets.

\newcommand{\folder}[1]{\gdef\grothFolder{#1}}
\newcommand{\batch}[1]{\gdef\grothBatch{#1}}
\newcommand{\leaves}[2]{\gdef\grothFirst{#1}\gdef\grothLast{#2}}
\newcommand{\foldertitle}[1]{\gdef\grothFoldertitle{#1}}
\gdef\grothFolder{?}\gdef\grothBatch{?}\gdef\grothFirst{?}\gdef\grothLast{?}\gdef\grothFoldertitle{}

% --- Le résumé --------------------------------------------------------------
%
% Il ouvre la lecture modernisée, sous le titre « Résumé », et il est composé
% en petit. Ce n'est pas un ornement : c'est le seul passage du document qui
% ne soit pas des mathématiques, et un lecteur doit pouvoir le reconnaître
% avant de l'avoir lu — puis le sauter s'il connaît déjà le sujet, ou s'y
% arrêter s'il ne le connaît pas. Le filet qui le suit marque la reprise du
% corps.

\newenvironment{resume}
  {\small\setlength{\parskip}{0.45em}}
  {\par\medskip\hrule\medskip}

% --- Mots-clés ---------------------------------------------------------------
%
% En anglais, en clôture du résumé de la lecture modernisée : le vocabulaire
% moderne sous lequel chercher ce que les feuillets construisent. Le manifeste
% du site (`npm run manifest`) les extrait pour étiqueter la cote entière —
% cette ligne est la source unique des tags, il n'y a pas de fichier à côté.
\newcommand{\keywords}[1]{\par\smallskip\noindent
  {\footnotesize\textsc{Keywords} --- \textit{#1}}\par}

% --- L'appareil critique ----------------------------------------------------

% Le numéro de feuillet, en marge. C'est l'ancre qui permet de régler un
% désaccord sur une formule en regardant *un* feuillet plutôt qu'une chemise
% de six cents. Le site s'en sert aussi pour tourner les pages du fac-similé
% quand on fait défiler la transcription.
\newcommand{\leaf}[1]{%
  \marginnote{\footnotesize\color{gray!70}\textsf{#1}}%
  \label{leaf:#1}\ignorespaces}

% Illisible : on ne devine pas. Un mot inventé qui se lit comme les autres est
% le pire résultat possible.
\newcommand{\ill}{\textcolor{red!60!black}{[\,\ldots\,]}}

% Lecture douteuse : le mot est donné, et signalé comme incertain.
\newcommand{\uncertain}[1]{\uline{#1}}

% Ajout de l'éditeur — une lettre manquante, un mot sous-entendu.
\newcommand{\add}[1]{\textcolor{blue!50!black}{[#1]}}

% Biffé par Grothendieck lui-même. Ce qu'il a rayé fait partie du document :
% une rature dit souvent où la pensée a changé de direction.
\newcommand{\struck}[1]{\sout{#1}}

% Note du transcripteur, sur l'état matériel du feuillet.
\newcommand{\note}[1]{\par\smallskip{\small\color{gray!60!black}\textit{[#1]}}\par\smallskip}

% Note marginale de Grothendieck, distincte du corps du texte.
\newcommand{\marginal}[1]{\par\smallskip\noindent\hspace{1em}%
  {\small\color{gray!60!black}#1}\par\smallskip}

% --- Titre ------------------------------------------------------------------

\AtBeginDocument{%
  \begin{center}
    {\large\bfseries Fonds Alexandre Grothendieck --- cote n\textsuperscript{o}~\grothFolder}\\[2pt]
    {\itshape\grothFoldertitle}\\[4pt]
    {\small feuillets \grothFirst--\grothLast\ (lot \grothBatch)}\\[2pt]
    {\footnotesize\color{gray!60!black}Transcription établie d'après le fac-similé numérisé
     par l'Université de Montpellier.\\
     Les lectures incertaines sont soulignées, les illisibles notées [\,\ldots\,],
     les ajouts entre crochets.}
  \end{center}
  \vspace{1em}}

\endinput


\folder{6}
\batch{1}
\pages{1}{20}
\dating{1966-[à partir de 1970]}
\watermark{Édition de démonstration}
\foldertitle{Lettre Tate (mai 66) (Cristaux) : lettre (1966), tapuscrit (s.d.), notes manuscrites (s.d.)}

\begin{document}

\page{1}
\marginal{Pise, Mai 1966}\note{à l'encre, en haut à droite, de sa main. La lettre est dactylographiée ; ses ajouts et corrections à l'encre (ou tapés dans l'interligne et appelés à l'encre) sont marqués comme ajouts, ses suppressions comme biffures, avec une note quand il le faut. Les lettres soulignées à la machine ($\underline{F}$, $\underline{Z}$, $\underline{M}$\ldots) restent soulignées ; les mots soulignés sont en italique. La lettre est paginée en haut à droite : 1 ici, puis 2 à la page 3 des archivistes, et ainsi de suite (page de la lettre = page des archivistes moins un).}

Cher John,

J'ai réfléchi aux groupes formels et à la cohomologie de De Rahm, et suis
arrivé à un projet de théorie, ou plutôt \add{de} début de théorie, que j'ai envie de
t'exposer\add{,} pour me clarifier les idées.

\section*{\emph{Chapitre} 1\quad \emph{La notion de cristal.}}

Commentaire terminologique : Un cristal possède deux propriétés caractéristiques :
la \emph{rigidité}, et la faculté de \emph{croitre}, dans un voisinage approprié.
Il y a des cristaux de toute espèce de substance : des cristaux de soude, de
souf\struck{f}re, de modules, d'anneaux, de schémas relatifs etc.

\add{1.0.}\note{le numéro « 1.0. » est ajouté à l'encre, dans la marge de gauche.}
Soit $S$ un préschéma, au dessus d'un autre $R$ ; dans le cas qui nous
intéressera le plus, on aura $R=\operatorname{Spec}(\underline{Z})$. Soit $C$ la
catégorie des \add{$R$-}préschémas \add{$T$} sous $S$, \add{(}i.e. munis d'un
$R$-morphisme $S\to T$,\add{)} tels que $S\to T$ soit une immersion fermée
définie par un idéal localement nilpotent.\note{« $R$- » et « $T$ » sont ajoutés
à l'encre au-dessus de la ligne, les parenthèses autour de « i.e.\ \ldots\ $S\to T$, » aussi.}
On a le foncteur « oubli » de $C$ dans $(\mathrm{Sch})$, et la catégorie fibrée
des Modules quasi-cohérents sur des préschémas variables induit donc une
catégorie fibrée sur $C$, associant à tout $T$ la catégorie des Modules
quasi-cohérents sur $T$.

\emph{Définition} 1.1. On appelle cristal de modules (sous-entendu :
quasi-cohérents) sur $S$, relativement à $R$, une section cartésienne de la
catégorie fibrée précédente au dessus de $C$. Plus généralement, pour toute
catégorie fibrée $\underline{F}$ sur $(\mathrm{Sch})_{/R}$, on définit la
notion de « $\underline{F}$-cristal » sur $S$, ou « cristal en objets de
$\underline{F}$ », de la façon correspondante. Ceci donne un sens aux
expressions : cristal en algèbres, en algèbres commutatives, en préschémas
relatifs etc\add{,} sur $S$ relativement à $R$. Quand
$R=\operatorname{Spec}(\underline{Z})$, on parlera de « cristal absolu » sur
$S$, de l'espèce considérée.

1.2. Les $\underline{F}$-cristaux sur $S$ forment une catégorie, qui dépend
fonctoriellement de $\underline{F}$. Ceci permet, en particulier, de définir sur
la catégorie des cristaux de modules sur $S$ les opérations tensorielles
habituelles : produit tensoriel de deux cristaux de modules etc, satisfaisant
aux propriétés habituelles.

\page{3}
\note{p.~2 de la lettre.}
1.3. Regardons de même ce qui se passe quand on fixe $\underline{F}$ et fait
varier $R,S$. Tout d'abord, si $S\to R'\to R$, alors tout cristal sur $S$
relativement à $R$ en définit un relativement à $R'$, par simple restriction.
En particulier, un cristal absolu définit un cristal relatif, pour tout
préschéma de base $R$ de $S$.

Fixons maintenant $R,S$ et soit $S'\to S$ un morphisme. On définit alors de
façon naturelle un foncteur « image inverse » allant des cristaux de type
$\underline{F}$ sur $S$ vers les cristaux de même type sur $S'$ (tout relatif
à $R$). Pour s'en assurer, il suffit de définir un foncteur $C'\to C$ (où $C'$
est défini en termes de $S'$ comme $C$ en termes de $S$), compatible avec les
foncteurs « oubli ». Or si $S'\to T'$ est un objet de $C'$, il y a une
construction évidente de somme amalgamée dans la catégorie des préschémas, qui
donne un $T=S\amalg_{S'}T'$ et un morphisme $S\to T$, qui fait de $T$ un objet
de $C$, ce qui définit le foncteur cherché.

On peut donc dire que pour un $S$ variable sur $R$, les cristaux de type
$\underline{F}$ sur $S$ forment une \emph{catégorie fibrée} sur
$(\mathrm{Sch})_{/R}$, grâce à la notion d'image inverse précédente.

Les deux variances (en $R$, et en $S$) sont compatibles dans un sens évident,
et peuvent être régardées comme provenant d'une variance en $(R,S)$
directement.

1.4. On a un foncteur naturel et évident qui va des cristaux de modules
(disons) sur $S$ (rel à $R$) dans des Modules quasi-cohérents sur $S$ : c'est
le foncteur « valeur en $S$ ». On fera attention que ce foncteur n'est en
général pas même fidèle (cf exemple \add{1.5.} plus bas).\note{« 1.5. » est
écrit à l'encre dans un blanc laissé par la machine.} Il est donc un peu
dangereux de vouloir considérer un cristal de modules sur $S$ comme étant un
Module quasi-cohérent $\underline{M}$ sur $S$, \emph{muni} d'une structure
supplémentaire, sa « structure cristalline ». Dans certains cas cependant
\add{(cf 1.8.)}, le foncteur cristaux de modules $\rightsquigarrow$ Modules est
fidèle et le point de vue précédent devient plus raisonnable,

1.5. \emph{Exemple} 1 : \emph{relations avec les vecteurs de Witt}. Supposons
que $S$ soit le spectre d'un corps \emph{parfait} $k$, et
$R=\operatorname{Spec}(\underline{Z})$. Soit $W$ l'anneau de Cohen de corps
résiduel $k$ : si $\mathrm{car}\,k>0$, c'est l'anneau des vecteurs de Witt
défini par $k$, si $\mathrm{car}.k=0$, c'est $k$ lui-même\add{;} \add{dans ce
dernier cas, supposons $k$ \emph{algébrique} sur $\mathbf{Q}$.}\note{ajout à
l'encre, écrit dans l'interligne et relié par un trait au point, qu'il
change en point-virgule ; « algébrique » est souligné.} Alors il est bien connu
que la catégorie $C$ de 1.0.

\page{4}
\note{p.~3 de la lettre.}
est équivalente à celle des $W$-algèbres locales, annulées par une puissance
de l'idéal maximal de $W$, à extension résiduelle triviale. Ecrivant
$W=\varprojlim W_n$ comme à l'accoutumé, dans le cas $p>0$, on trouve que la
donnée d'un cristal de modules (absolu) sur $k$ équivaut à celle d'un système
projectif $(M_n)$ « $p$-adique » de modules $M_n$ sur les $W_n$ ; donc les
cristaux de modules de présentation finie sur $k$ forment une catégorie
équivalente à celle des modules de type fini sur $W$. Si $p=0$, alors la
catégorie des cristaux de modules sur $k$ est équivalente à celle des
vectoriels sur $k=W$. Ces descriptions sont compatibles avec les notions de
changement de base pour $k$ variable. Elles se formulent évidemment pour toute
autre espèce de cristaux, défini par une catégorie fibrée $\underline{F}$.

Dans le cas envisagé, le foncteur \add{« }valeur en $S$\add{ »}, sur la
catégorie des cristaux de présentation finie sur $k$, s'identifie au foncteur
$\otimes_W k$ sur la catégorie des modules de type fini sur $W$, foncteur qui
(si $p>0$) n'est pas fidèle.

1.6. \emph{Exemples} 2 : \add{$S$ étale sur $R$.}\note{ajout à l'encre
au-dessus de « Si $S=R$ », que le trait d'appel embrasse.} Si $S=R$, la
catégorie $C$ admet $R$ lui-même comme objet final, donc le foncteur « valeur
en $R$ » est une équivalence de la catégorie des cristaux sur $R$, de type
$\underline{F}$ donné, avec la catégorie $\underline{F}_R$. En particulier, un
cristal de modules sur $\operatorname{Spec}\underline{Z}$ est essentiellement
la même chose qu'un $\underline{Z}$-module.

De façon un peu plus générale, si $S$ est étale sur $R$, alors $S$ est un
objet final de $C$, et les cristaux de modules (disons) sur $S$, relativement à
$R$, s'identifient aux Modules quasi-cohérents sur $S$.

1.7. \struck{Relations avec la notion d'objet stratifié} \emph{Exemple} 3 :
$S$ \emph{un sous-}\struck{objet de} \emph{préschéma de} $R$.\note{les mots
biffés le sont à la machine.} Comme la notion de cristal relatif sur $S$ ne
change pas si on remplace $R$ par un ouvert par lequel se factorise $S$, on peut
supposer $S$ fermé dans $R$, défini par un Idéal quasi-cohérent $J$. Soit
$S_n=V(J^{n+1})$ le $n$.ème voisinage infinitésimal de $R$ dans
$S$.\note{sic : lire « de $S$ dans $R$ ».} Alors la famille des objets $S_n$
de $C$ est finale dans $C$, d'où on conclut facilement qu'un cristal de type
$\underline{F}$ sur $S$ s'identifie à une suite $(M_n)_{n\in\underline{N}}$
d'objets des $\underline{F}_{S_n}$ qui se recollent. En particulier, si $R$ est
localement noethérien, un cristal de modules sur $S$ relativement à $R$, de
présentation finie, s'identifie à un Module cohérent sur le complété formel de
$R$ le long de $S$.

\page{5}
\note{p.~4 de la lettre.}
Cet exemple contient le cas de $\mathrm{car}.p>0$ de l'éxemple 1, si on note
qu'a priori, grâce à Cohen, un cristal absolu sur $k$\add{,} c'est pareil
qu'un cristal relativement à $W$.

On peut aussi donner un énoncé commun aux exemples 2 et 3, en partant du cas
où on se donne un $S\to R$ non \emph{ramifié}, ce qui permet en effet de
construire encore des « voisinages infinitésimaux » $S_n$.

1.8. \emph{Relation avec la notion de stratification}. Les données
$R,S,\underline{F}$ étant comme d'habitude, considérons pour chaque entier
$n\geqslant 0$ le voisinage infinitésimal $\Delta_n$ de la diagonale de
$S\times_R S$, qui s'envoie dans $S$ par les deux projections
$\mathrm{pr}_1$ et $\mathrm{pr}_2$. Si $E$ est un objet de
$\underline{F}_S$, une $n$-\emph{connexion} sur $E$ (relativement à $R$) est
la donnée d'un isomorphisme $\mathrm{pr}_1^*(E)\simeq\mathrm{pr}_2^*(E)$ qui
induit l'identité sur la diagonale. Une $\infty$-\emph{connexion} ou
\emph{pseudo-stratification} de $E$, est la donnée pour tout $n$ d'une
$n$-connexion, de telle façon que ces $n$-connexions se recollent. Enfin, une
\emph{stratification} sur $E$ est la donnée d'une pseudo-stratification qui
satisfait aux conditions formelles d'une donnée de descente, quand on fait
intervenir les voisinages infinitésimaux de tous ordres de \add{la diagonale
de} $S\times_R S\times_R S$.\note{« la diagonale de » est ajouté à l'encre
au-dessus de la ligne.} Ces notions donnent lieu à des sorites analogues à
ceux de 1.2.\ et 1.3. Notons que lorsque $S$ est « formellement non ramifié
sur $R$ » i.e.\ $\underline{\Omega}^1_{S/R}=0$, alors toutes ces notions
deviennent triviales : un objet $E$ admet alors toujours exactement une
connexion de tout ordre, donc une seule pseudo-stratification, et celle-ci est
une stratification : donc la catégorie des objets de $F_S$ \emph{munis} d'une
stratification relativement à $R$ est alors équivalente, par le foncteur
« valeur en $S$ », à la catégorie $\underline{F}_R$ elle même.\note{sic
($\underline{F}_R$ pour $\underline{F}_S$ ?).} (Dans tous les cas, le foncteur
« valeur en $S$ » est \struck{en tous cas} fidèle).\note{« en tous cas » biffé
à l'encre.}

Ceci \uncertain{défini},\note{correction à l'encre sur le mot tapé, qu'elle
recouvre.} on voit que pour tout cristal $\underline{M}$ sur $S$ de type
$\underline{F}$ (relativement à $R$), sa valeur
$\underline{M}(S)=M$ est un objet de $\underline{F}_S$ muni d'une
\emph{stratification canonique} relativement à $R$, d'où un foncteur :
cristaux relatifs de type $\underline{F}$ $\rightsquigarrow$ objets de
$\underline{F}$ munis d'une stratification. La remarque de 1.4.\ montre
d'ailleurs que ce foncteur n'est pas toujours fidèle. Il y a cependant des cas
intéressants

\page{6}
\note{p.~5 de la lettre.}
où ce foncteur est une \emph{équivalence de catégories} : il en est en tous cas
ainsi si $S\to R$ est « \emph{formellement lisse} », par exemple si c'est un
morphisme lisse, ou si $S$ et $R$ sont des spectres de corps $k_0,k$, avec $k$
une extension \emph{séparable} de $k_0$. \struck{Dans ce dernier cas}(Quand
d'ailleurs $S\to R$ est même formellement étale, alors les deux catégories :
celle des cristaux et celle des objets à stratification, deviennent
\struck{d'ailleurs} équivalentes\add{,} par le foncteur « valeur en
$S$ »\add{,} à $\underline{F}_S$ lui-même, ce qui nous redonne l'exemple à la
noix de 1.5 où $k$ est de car.\ nulle).\note{« Dans ce dernier cas » est
biffé à la machine, « d'ailleurs » à l'encre.} Sauf erreur, on a encore une
équivalence de catégories si $S\to R$ est \emph{plat et localement de
présentation finie}, \add{(du moins si $R$ loc.\ noeth., et se bornant aux
Modules cohérents),} mais je n'ai pas écrit de démonstration.\note{ajout à
l'encre dans l'interligne, relié par un trait à la virgule ; le soulignement
de « plat et localement de présentation finie » est de sa main.}

1.9. \emph{Une digression sur la notion de stratification en caractéristique
nulle.} Quand $S$ est lisse sur $R$ (en fait, il suffit que $S$ soit
différentiellement lisse sur $R$, i.e.\ le morphisme diagonal
$S\to S\times_R S$ une immersion régulière), et si $R$ est de caractéristique
nulle, alors une stratification d'un Module $M$ sur $S$ (relativement à $R$)
est connue quand on connait la $1$-connexion qu'elle définit, et on peut se
donner celle-ci arbitrairement, soumise à la seule condition que le « tenseur
courbure », qui est une certaine section de
$\underline{\Omega}^2_{S/R}\otimes\underline{\mathrm{End}}(M)$, soit nul.
(Cela peut aussi s'exprimer en disant qu'on fait opérer le faisceau
$\underline{\mathrm{Dér}}_{S/R}$ des dérivations relatives de $S$ sur $R$, sur
$M$, de façon à satisfaire aux relations habituelles, y compris celle de
transformer crochet en crochet). J'ignore dans quelle mesure
\struck{la propriété}\add{l'hypothèse} de lissité est nécessaire pour cet
énoncé,\note{« propriété » biffé, « hypothèse » tapé dans l'interligne,
« la » corrigé en « l' ».} ou si (dans le cas lisse disons) on peut formuler
un énoncé analogue pour toute catégorie fibrée $\underline{F}$ sur
$(\mathrm{Sch})_{/R}$. Mais bien entendu, l'hypothèse de caractéristique nulle
faite sur $R$ est tout à fait essentielle. Si $R$ est de caractéristique $p>0$,
l'énoncé qui remplace le précédent (et qui sauf erreur est dû à Cartier) est
que dans ce cas, la donnée d'une « connexion sans torsion » sur $M$
\add{« compatible avec puissances $p$.ièmes »}\note{ajout à l'encre dans la
marge de gauche, appelé par un trait après « $M$ ».} équivaut à une « donnée
de descente » sur $M$ relativement à frobenius $S\to S^{(p/R)}$. Il y a loin
de là à une \struck{rigidifi} stratification !

1.10. \emph{La notion de $p$-cristal et ses variantes.}

Nous supposons maintenant que $S$ est de caractéristique $p>0$,
\struck{\ill{}}\note{un mot biffé à la machine en fin de ligne.}

\page{7}
\note{p.~6 de la lettre.}
et $R=\operatorname{Spec}(\underline{Z})$ (ou
$\operatorname{Spec}(\underline{Z}_p)$, spectre des entiers $p$-adiques, cela
reviendrait au même). Si $\underline{M}$ est un cristal de modules sur $S$,
\add{de présentation finie,} alors en vertu de 1.5., pour tout $s\in S$, si
$s'$ est le spectre d'une clôture parfaite de $k(s)$, l'image inverse
$\underline{M}(s')$ de $\underline{M}$ en $s'$ peut être interprétée comme un
module de type fini sur l'anneau des vecteurs de Witt $W(k(s'))$.\note{« de
présentation finie, » est tapé dans l'interligne et appelé à l'encre.} Ainsi,
la notion de cristal de modules sur $S$ (au sens absolu) semble convenable pour
formaliser la notion de \emph{« famille \add{algébrique »} de modules sur les
anneaux de vecteurs de Witt aux différents points parfaits au dessus de}
$S$.\note{« algébrique" » est tapé dans l'interligne et appelé à l'encre ; le
soulignement est en partie de sa main.} Bien entendu, la notion de cristal est
plus fine que ça encore, et doit donner, j'espère, la bonne notion de
« famille » lorsque, disons, $S$ est le spectre d'un corps de fonctions (donc
pas nécessairement parfait). Je vais maintenant introduire une structure
supplémentaire, qui devrait permettre de formuler, de même, la notion de
« famille algébrique de modules de Dieudonné », paramétrée par $S$,

Considérons le morphisme « puissance $p$.ème »
\[
  S\xrightarrow{\ \mathrm{frob}_S\ }S\ ,
\]
il permet d'associer, à tout cristal (absolu) sur $S$, d'espèce quelconque, un
cristal image inverse :
\[
  \underline{M}^{(p)}=\mathrm{frob}_S(\underline{M})\ .
\]
On appelle $p$-\emph{cristal} sur $S$ un cristal $\underline{M}$ sur $S$, muni
d'un morphisme de cristaux
\[
  F\colon \underline{M}^{(p)}\longrightarrow\underline{M}\ .
\]
Evidemment, les $p$-cristaux \add{sur $S$} d'espèce $\underline{F}$ donnée
forment encore une catégorie, dépendant fonctoriellement de $\underline{F}$
(pour les foncteurs \emph{covariants} cartésiens de catégories fibrées), ce qui
permet par exemple, pour les $p$-cristaux de modules sur $S$, d'introduire
toutes les opérations tensorielles habituelles covariantes (produits
tensoriels, puissances alternées et symétriques etc). Pour \struck{formuler}
définir le \emph{dual} d'un $p$-cristal de modules, il y a lieu d'introduire
une notions duale de celle de $p$-cristal d'espèce $\underline{F}$, c'est celle
de $p^{-1}$-\emph{cristal} d'espèce $\underline{F}$ : c'est un cristal
d'espèce $\underline{F}$, avec un morphisme de cristaux en sens inverse :

\page{8}
\note{p.~7 de la lettre.}
\[
  V\colon \underline{M}\longrightarrow\underline{M}^{(p)}\ .
\]
Ainsi un \add{contra-}foncteur cartésien entre catégories fibrées sur
$(\mathrm{Sch})$ transforme $p$-cristaux en $p^{-1}$-cristaux, et inversement.
Les $p^{-1}$-cristaux de modules se multiplient tensoriellement entre eux comme
les $p$-cristaux de modules.

Pour obtenir une notion « autoduale », il y a lieu d'introduire la notion de
\emph{bi-cristal} \add{\emph{de poids} 1} sur $S$, qu'on pourrait aussi appeler
un \emph{cristal de Dieudonné} sur $S$, lorsque $\underline{F}$ est fibré en
catégories additives : c'est un cristal muni à la fois d'une $p$-structure et
d'une $p^{-1}$-structure, satisfaisant aux relations
\[
  FV=p.\,\mathrm{id}_M,\qquad VF=\mathrm{id}_{M^{(p)}}\ .
\]
\note{« contra- » et « de poids 1 » sont tapés dans l'interligne et appelés à
l'encre ; le « $p.$ » manque dans la seconde relation, qui est donnée comme
tapée.}

Cette notion semble tout à fait adéquate à l'étude des groupes formels, ou ce
qui revient moralement au même, à l'étude de la cohomologie de De Rham en
dimension 1. En dimension\struck{s} supérieure\struck{s} \add{$i$}, il y a
lieu d'introduire la notion de \emph{bi-cristal de poids} $i$, qui est un
cristal muni de $F$ et $V$ satisfaisant aux relations
\[
  FV=p^i\,\mathrm{id}_M\ ,\qquad VF=p^i\,\mathrm{id}_{M(p)}\ .
\]
Par exemple, on définit le \emph{bicristal de Tate de poids} $2i$, qui est
défini par le cristal de modules unité (associant à tout $T$ sous $S$ le
Module $\underline{O}_T$ lui-même), et les morphismes de cristaux
\[
  F=p^i\,\mathrm{id}_{T^i}\ ,\qquad V=p^i\,\mathrm{id}_{T^i}\ .
\]
\marginal{$2i$}\note{dans la marge de gauche, entouré, à hauteur de « poids
$2i$ », où le chiffre tapé a été refrappé ; l'indice est tapé « $T\,i$ ».}
Ainsi le bicristal de Tate de poids $2i$ est la puissance tensorielle $i$.ème
du bicristal de Tate de poids 1.\note{sic : « poids 1 » ; d'après ce qui
précède on attendrait le poids 2, celui de $T^1$ en 1.11.} (NB les bicristaux
de poids quelconques se multiplient tensoriellement, de \struck{sorte}
\add{façon} que les poids s'ajoutent).\note{« façon » à l'encre, sur le mot
tapé.}

1.11. Si on veut « rendre inversible » le bicristal de Tate $T^1$, on est
obligé, qu'on le veuille ou non, de passer à une catégorie de fractions de la
catégorie des cristaux de modules sur $S$, [qui entre parenthèses est une
catégorie $\underline{Z}_p$-linéaire, (i.e.\ les Hom sont des
$\underline{Z}_p$-modules \ldots), tout comme les catégories de $p$-cristaux
etc qu'on en déduit]. Il revient donc au même de passer à une catégorie de
fractions, en déclarant qu'on veut rendre notre catégorie abélienne
$\underline{Q}$-linaire, ou en déclarant qu'on la veut rendre
$\underline{Q}_p$-linéaire : il faut faire le quotient par la sous-catégorie
abélienne épaisse formée des objets de torsion, qui sont aussi des objets de
$p$-torsion \add{(et cela revient simplement à garder les mêmes objets, et à
prendre comme nouveaux Hom les $\mathrm{Hom}\otimes_{\underline{Z}}\mathbf{Q}$.)}.\note{ajout
à l'encre au pied de la page, encadré et appelé après « $p$-torsion » ; les
crochets sont aussi de sa main.} On trouve la catégorie des \emph{« cristaux
de modules}

\page{9}
\note{p.~8 de la lettre.}
\emph{à isogénie près »}, ou \emph{isogristaux}\note{sic.}, sur $S$. Utilisant cette
nouvelle notion, on définit comme ci-dessus la notion de $p$-isocristal, comme
étant la donnée d'un isocristal $\underline{M}$ muni d'un homomorphisme
$F\colon\underline{M}^{(p)}\to\underline{M}$. La notion de bi-isocristal de
poids $i$, qui serait calquée de celle de bicristal de poids $i$, n'est pas
très raisonnable alors, car $V$ doit être alors donné en termes de $F$ comme
$p^iF^{-1}$. Il y a lieu plutôt d'appeler \emph{bi-isocristal} (sans
précision de poids) un $p$-isocristal pour lequel $F$ est un
\emph{isomorphisme}, et de ne pas choisir entre les différents $p^iF^{-1}$
possibles ; il vaut mieux également ne pas parler ici de
$p^{-1}$-isocristal, la notion intéressante étant celle de bi-isocristal, qui
est manifestement \emph{autoduale} \struck{: si $\phi$ est un foncteur
contravariant de catégories fibrées sur} \add{quand on se restreint aux objets
de type fini :} le dual de $(\underline{M},F)$ est
$(\underline{M},{}^tF^{-1})$, où $\underline{M}$ est le iso-cristal dual de
$M$.\note{la ligne biffée l'est à la machine et à l'encre ; « quand on se
restreint aux objets de type fini : » est tapé dans l'interligne. Aucun signe
ne distingue, dans « $(\underline{M},{}^tF^{-1})$ », le dual de
$\underline{M}$ : sic.}

On peut également, bien sûr, localiser par rapport à $p$ en partant déjà de la
catégorie des bi-cristaux de poids $i$, on trouve les \emph{bicristaux de
poids} $i$ \emph{à isogénie près}, qui forment une catégorie abélienne
$\mathrm{Isbicr}(S,i)$, et un foncteur exact « oubli de $V$ »
\[
  \mathrm{Isbicr}(S,i)\longrightarrow\mathrm{Isbicr}(S)\ ,
\]
à valeurs dans la catégorie $\mathrm{Isbicr}(S)$ des iso-bicristaux sur $S$.
\note{dans les deux premières occurrences, « bi » est ajouté au sigle
« Iscr » tapé d'abord.} Ce foncteur est \struck{\ill{}} \emph{pleinement
fidèle}, de sorte que \emph{les \uncertain{bi-isocristaux} forment une
généralisation commune des bicristaux de poids} $i$ \emph{à isogénie près}.
\note{un mot à l'encre, biffé, précède « pleinement fidèle » ; le préfixe de
« \ldots cristaux » est refrappé et se lit mal.} Pour préciser ce point,
désignant par $\mathrm{Bicr}(S,i)$ la catégorie des bicristaux de poids $i$ sur
$S$, il y a lieu d'introduire des foncteurs canoniques
\[
  \mathrm{Bicr}(S,i)\longrightarrow\mathrm{Bicr}(S,i+1)\longrightarrow\cdots,
\]
donnés par $(\underline{M},F,V)\rightsquigarrow(\underline{M},F,pV)$. Quand on
localise ces foncteurs par $p$, on trouve des foncteurs
\[
  \mathrm{IsBicr}(S,i)\longrightarrow\mathrm{IsBicr}(S,i+1)\longrightarrow\cdots
\]
qui se trouvent être \emph{pleinement fidèles} : on peut donc considérer que,
pour $i$ croissant, la notion de « bicristal de poids $i$, à isogénie près »
constitue une généralisation de plus en plus large de celle de « cristal de
Dieudonné ». La noti-

\page{10}
\note{p.~9 de la lettre.}
on de « bi-isocristal » est à ce titre une \struck{\ill{}}
\add{généralisation} qui englobe toutes les précédentes, et qui de plus a la
louable vertu d'être stable par produit tensoriel\add{ et passage au dual.}
Cela semble donc une catégorie toute indiquée comme catégorie des valeurs pour
un foncteur « cohomologie de De Rham » convenable, lorsqu'on se décide à
travailler à isogénie près.\note{« généralisation » est tapé au-dessus du mot
biffé à la machine ; « et passage au dual. » est ajouté à l'encre, la virgule
qui suivait « tensoriel » changée en point-virgule par la même main.}

1.12. \emph{Exemple} 1 : \emph{Cas où $S$ est le spectre d'un corps parfait}
$k$. Alors la donnée d'un cristal de modules de type fini sur $k$ équivaut à la
donnée d'un module de type fini $M$ sur $W$\add{,} la formation du motif
$\underline{M}^{(p)}$ correspond à la formation de
\[
  M^{(p)}=M\otimes_W(W,f_W)\ ,
\]
où $f_W$ est l'endomorphisme de frobenius de $W$, et $(W,f_W)$ est la
$W$-algèbre définie par $f_W$. Par suite un structure de $p$-cristal sur $M$
revient à la donnée d'un homomorphisme de $W$-modules $M^{(p)}\to M$, ou si on
préfère, à la donnée d'un homomorphisme $f_W$-semi-linéaire
\[
  F_M\colon M\longrightarrow M\ .
\]
Comme l'application $x\mapsto x\otimes 1$ de $M$ dans $M\otimes_W(W,f_W)$ est
bijective, $f_W$ étant un automorphisme de $W$, on peut considérer la bijection
inverse, qui est $f_W^{-1}$-semi-linéaire. Par suite, la donnée d'une
$p^{-1}$-structure sur $M$ revient à la donnée d'un homomorphisme
$f_W^{-1}$-linéaire :
\[
  V_M\colon M\longrightarrow M\ .
\]
La donnée d'un couple $(F,V)$ définit sur $M$ une structure de bi-cristal de
poids $i$ si et seulement si on a
\[
  FV=VF=p^i\,\mathrm{id}_M\ .
\]
En particulier, pour $i=1$, on retrouve la notion de \emph{module de
Dieudonné} : la catégorie des cristaux de Dieudonné sur $k$ est canoniquement
équivalente à celle des modules de Dieudonné relativement à $k$.

La catégorie des isocristaux de type fini sur $k$ est équivalente à la
catégorie des vectoriels de dimension finie sur le corps des fractions
\add{$K$} \struck{\ill{}} de $W$. Donc un bi\add{iso}cristal (de type fini) sur
$k$ s'identifie à un tel vectoriel, muni d'un $f_K$-endomorphisme
bijectif.\note{« $K$ » est ajouté à l'encre en fin de ligne, le mot qui
commence la ligne suivante biffé à l'encre ; « iso » est tapé dans
l'interligne et appelé à l'encre ; un exposant tapé après $f_K$ est biffé.}

\page{11}
\note{p.~10 de la lettre.}
Toutes ces descriptions sont compatibles avec les opérations tensorielles, et
la formation d'images inverses pour $k$ variable.

On peut se demander, avec la description précédente des bi-isocristaux, quels
sont ceux qui sont dans l'image essentielle de
$\mathrm{Isbicr}(S,i)\to\mathrm{Isbicr}(S)$, On trouve que ce sont les couples
$(E,F)$, $E$ un vectoriel sur $K$ et $F$ un automorphisme semi-linéaire, tels
que, si on pose $V=p^iF^{-1}$, pour tout $x\in E$, l'ensemble de ses
transformés par les différents monômes non commutatifs en $F,V$ soit une
partie \emph{bornée} de $E$ : c'est une pure tautologie.\note{la lettre $E$
est refrappée sur un $V$.} Pour qu'il existe un $i$ ayant cette propriété, il
faut et il suffit que pour tout $x\in E$, l'ensemble des $F^nx$ soit borné.
Bien entendu, c'est là une propriété qui est stable par changement de $F$ en
$pF$ (correspondant à la tensorisation par le bi-isocristal de Tate de poids
2), mais non par le changement de $F$ en $p^{-1}F$ (correspondant à la
tensorisation par l'inverse $T^{-1}$ du bi-isocristal précédent). Comme on
tient beaucoup à inverser Tate, on voit qu'il n'est pas naturel, en effet, de
poser une condition de croissance sur l'ensemble des itérés de $F$. De façon
précise, appelant \emph{effectifs} les bi-isocristaux qui proviennent d'un
objet d'un $\mathrm{Isbicr}(S,i)$, on voit que tout bi-isocristal est de la
forme $T^{-i}\otimes\underline{N}$, avec $\underline{N}$ effectif : on n'a pas
ajouté plus de nouveaux objets qu'il n'était absolument nécessaire !

1.13. \emph{Exemple} 2 : $S$ \emph{lisse sur un corps parfait} $k$.
Déterminons d'abord dans ce cas les cristaux sur $S$, sans plus. Tout d'abord,
sans condition de lissité, on voit à l'aide de la propriété caractéristique de
$W$ que la catégorie $C$ introduite dans 1.0.\ ne change pas\add{,}
essentiellement\add{,} quand on remplace $R=\operatorname{Spec}(\underline{Z})$
par $R=\operatorname{Spec}(W)$. D'autre part, il résulte aussitôt des
définitions que la donnée d'un cristal sur $S$, relativement à $W$, revient à
la donnée d'un système cohérent de cristaux sur $S$, relatifs aux
$W_n=W/p^{n+1}W$. (NB On n'a pas formulé avec la généralité qui convenait
l'exemple 3 de 1.7. ! ) Théoriquement, on est donc ramené à déterminer, pour
chaque $n$, la catégorie des cristaux sur $S$ relatifs à $W_n$, et les
foncteurs restriction entre ces catégories. D'ailleurs, il est inoffensif de se
localiser sur $S$, par exemple de supposer au besoin $S$ affine. Utilisant
maintenant la lissité de $S$, on peut donc supposer que pour tout $n$, $S$ se
remonte en

\page{12}
\note{p.~11 de la lettre.}
un $S_n$ lisse sur $W_n$, de façon que ces choix se recollent (il suffit pour
ceci $S$ lisse et affine). Or, il résulte encore facilement des définitions,
utilisant seulement que $S=S_0\to S_n$ est une immersion fermée définie par un
Idéal nilpotent, que le foncteur restriction des cristaux sur $S_n$
(relativement à une base quelconque -- ici on prendra $W_n$) vers les
cristaux sur $S_0$, est une équivalence de catégories, donc un cristal sur $S$
relativement à $W_n$ s'identifie à un cristal \add{sur $S_n$} relativement à
$W_n$. Comme $S_n$ est lisse sur $W_n$, il résulte de ce qui a été dit dans
1.8.\ que les cristaux en question s'identifient aux objets (de l'espèce
$\underline{F}$ considérée) sur $S_n$, \emph{munis d'une stratification
relativement à} $R_n=\operatorname{Spec}(W_n)$. Les foncteurs restrictions sur
les cristaux s'expriment alors par les foncteurs restrictions correspondants
pour objets stratifiés. En résumé : \emph{un cristal} \add{$\underline{M}$}
\emph{sur} $S$ \emph{s'identifie à un système cohérent} $(M_n)$ \emph{d'objets
à stratification sur les différents} $S_n$ \emph{sur}
$R_n=\operatorname{Spec}(W_n)$.\note{« sur $S_n$ » et « $\underline{M}$ »
sont tapés dans l'interligne et appelés à l'encre.} Pour relier ceci à des
objets qui soient plus dans la nature d'objets « définis en caractéristique
nulle », supposons d'abord que l'on puisse même relever $S$ en un schéma $X$
\emph{propre sur} $R=\operatorname{Spec}(W)$, hypothèse évidemment bien
restrictive. On voit alors, utilisant les théorèmes de comparaison EGA III 4,5
que \emph{dans le cas de cristaux de modules cohérents sur} $S$, \emph{la
catégorie desdits est canoniquement équivalente à la catégorie des Modules
cohérents} $\underline{M}$ \emph{sur} $X$, \emph{munis d'une stratification
relativement à} $R=\operatorname{Spec}(W)$. (NB Il se trouve donc, à
posteriori, que cette dernière catégorie est essentiellement indépendante du
relèvement choisi $X$ de $S$). Considérant la fibre générique $X_K$ de $X$ sur
$R$, qui est un schéma propre et lisse sur un corps de caractéristique nulle,
on trouve donc un foncteur remarquable « restriction » ou « localisation »,
allant de la catégorie des cristaux de modules cohérents sur $S$, dans la
catégorie des Modules cohérents sur $X_K$, stratifiés relativement à $X_K$. Or
(oubli de 1.8.)\ on voit facilement qu'un Module cohérent stratifié sur un
préschéma localement de type fini sur un corps est nécessairement localement
libre. De plus comme ici $K$ est de caractéristique nulle\add{ et $X_K$ lisse
dessus}, on a signalé dans 1.8.\ que la notion de stratification s'explicite
très simplement comme celle de « connexion inté-\note{« et $X_K$ lisse
dessus » est tapé dans l'interligne et appelé à l'encre. « stratifiés
relativement à $X_K$ » : sic (on attendrait « à $K$ »).}

\page{13}
\note{p.~12 de la lettre.}
grable ». [Enfin, lorsque $S$ donc $X_K$ est géométriquement connexe, et que
$K$ peut se plonger dans le corps des complexes $\underline{C}$, alors la notion
de module cohérent à connexion intégrable sur $X_K$ s'interprète en termes de
représentations linéaires \add{(complexes)} du groupe fondamental transcendant
de $X_{\underline{C}}^{\mathrm{an}}$, de façon bien connue.\note{« ex » est
récrit à l'encre sur le mot tapé « connexion », mal frappé ; « (complexes) »
est tapé dans l'interligne, appelé à l'encre, avec un signe de vérification en
marge droite.} Si par exemple le groupe fondamental transcendant est le groupe
unité, alors on conclut par descente que tout Module cohérent stratifié sur
$X_K$ est trivial, ce qui en termes de $S$ s'énonce en disant que tout cristal
de modules cohérents sur $S$ est isogène à un cristal « constant », i.e.\
\add{à un cristal} qui est l'image inverse d'un cristal sur $k$, \add{(}lui-même
défini par un module de type fini sur $W$\add{)}.\note{« à un cristal » à
l'encre au-dessus d'un mot biffé à la machine ; les parenthèses sont de sa
main.} De ceci et de la rigidité de la notion de cristal on déduit\add{,} par
exemple\add{,} facilement que tout $p$-cristal cohérent sur $S$ ou tout
bi-cristal de poids $i$ donné (par exemple tout cristal de Dieudonné) est
isogène à un fourbis de même espèce \emph{trivial}.\note{« facilement » est
entouré à l'encre et les virgules qui l'isolent ajoutées, comme pour le
déplacer.} On voit donc là un principe d'approche transcendante pour l'étude
des familles de groupes formels (par exemple) en caractéristique $p$
\ldots]\note{les crochets sont tracés à l'encre.} Quand à la notion
d'isocristal (de modules cohérents) sur $S$, on constate aussitôt que le
foncteur précédent induit un foncteur \emph{pleinement fidèle} de la catégorie
de ces derniers, dans la catégorie des modules cohérents stratifiés sur $X_K$.
Il s'impose évidemment d'en déterminer l'image essentielle, et pour commencer
d'examiner \struck{quelle est l'im} si par hasard ce foncteur ne serait pas une
équivalence de catégories. Cela me semble un peu \emph{trop
optimiste}\marginal{Effectivement, ça \uncertain{ne marche} pas !}, mais je ne
suis sûr de rien, faute d'avoir regardé.\note{« trop optimiste » est souligné
d'un trait ondulé à l'encre qui mène à la note marginale, à gauche.} Tout ce
qu'on peut dire à priori, c'est que la condition cherchée sur un Module
stratifié doit être\struck{, à priori} de nature locale en les points de $X$
qui sont maximaux dans $S$.

Quand on ne suppose pas $S$ \struck{remonté} propre et remonté globalement,
mais qu'on suppose seulement qu'on a remonté $S$ formellement, en un
\emph{schéma formel} $X$ sur $R$, alors il est vrai (en fait, de façon
essentiellement triviale) que \emph{la catégorie des cristaux de modules
cohérents sur} $S$ \emph{est équivalente à la catégorie des Modules cohérents
sur le schéma formel} $X$, \emph{munis d'une stratification relativement à}
$\operatorname{Specf}(W)=\mathcal{R}$ (quand on transcrit de façon évidente
toutes les définitions envisagées dans 1.8.\ dans le cadre formel).\note{le
$\mathcal{R}$ final est de sa main.} Si on veut encore,

\page{14}
\note{p.~13 de la lettre.}
comme il est légitime, trouver un analogue pour la restriction à $X_K$
\add{envisagée} plus haut, il faut définir $X_K$ comme \emph{l'espace
rigide-analytique sur $K$ défini par le schéma formel} $X$, et considérer sur
$X_K$ des Modules cohérents munis de \struck{connexions} \emph{stratifications
au sens rigide-analytique}, ou ce qui revient au même, de connexions
intégrables \add{en ce sens}. De tels Modules sont encore nécessairement
localement libres.\note{« envisagée » est à l'encre ; « en ce sens » est tapé
dans l'interligne et appelé à l'encre.} On trouve ainsi un foncteur des
cristaux de Modules cohérents sur $S$ dans les Modules cohérents stratifiés sur
$X_K$, induisant un foncteur pleinement fidèle de la catégorie des isocristaux
cohérents sur $S$, dans la catégorie des Modules cohérents stratifiés sur
$X_K$. Comme tout à l'heure (et même plus, si on peut dire, car c'est vraiment
ici la situation « naturelle »), il s'impose de regarder quelle est l'image
essentielle. On peut également se demander si les Modules cohérents stratifiés
sur un espace rigide-analytique n'admettraient pas une description simple, en
termes d'un groupe plus ou moins discret jouant le rôle du groupe fondamental
dans la théorie transcendante sur le corps des complexes.

La description donnée \add{des cristaux sur $S$}, toute triviale qu'elle soit,
a déjà des conséquences intéressantes pour la structure de la catégorie des
cristaux \add{de modules} cohérents sur $S$ : cette catégorie est
noethérienne (si $S$ est noethérien i.e.\ de type fini sur $k$), tout objet
contient donc un plus grand sous-objet de torsion, de plus les objets sans
torsion correspond\add{ent} dans la description ci-dessus aux Modules
stratifiés sans torsion sur $X$, lesquels \add{sont} alors automatiquement
localement libres. Il est bien probable que des résultats analogues doivent
être vrais\add{,} sans hypothèse du genre lissité sur $S$.\note{les ajouts de
ce paragraphe sont tapés dans l'interligne et appelés à l'encre, sauf la
virgule, à l'encre ; un numéro tapé en tête de ligne, avant « La
description », est biffé à la machine.}

1.14. Il faut encore expliciter, dans la description générale précédente, le
\struck{donnée} foncteur $M\rightsquigarrow M^{(p)}$, pour pouvoir décrire de
façon analogue les bicristaux et bi-isocristaux. De façon générale, si en plus
de la donnée de $S$ sur $k$, on a un $S'$ lisse sur $k'$ parfait, et des
morphismes $S'\to S$ et $\operatorname{Spec}(k')\to\operatorname{Spec}(k)$
donnant lieu à un carré commutatif, et si enfin on peut relever $S'$
formellement en $X'$, et $S'\to S$ en $X'\to X$, donnant un carré commutatif
avec $\operatorname{Spec}W'\to\operatorname{Spec}W$, alors la notion d'image
inverse de cristaux de modules cohérents, de $S$ à $S'$, s'explicite en termes
d'image inverse de Modules cohérents stratifiés, de $X$

\page{15}
\note{p.~14 de la lettre.}
à $X'$. Lorsque $k=k'$, $S=S'$, et que les morphismes envisagés sont les
puissances $p$.èmes, on est donc conduit à chercher à \struck{prolonger}
relever ce morphisme à $X$ de façon compatible avec le morphisme
$\operatorname{Specf}(W)\to\operatorname{Specf}(W)$ déduit de $f_W$, ou ce qui
revient au même, de relever le morphisme canonique $S\to S^{(p/k)}$ en un
morphisme de $R$-schémas formels $X\to X\times_R(R,f_R)$.\note{le « f » de
« Specf » paraît récrit sur la parenthèse tapée.} C'est en tous cas toujours
possible si $S$ est affine, cas auquel on peut se ramener. Ayant donc ainsi un
morphisme
\[
  f_X\colon X\longrightarrow X\ ,
\]
compatible avec $f_R$ sur $R=\operatorname{Specf}(W)$, et induisant $f_S$ sur
$S$, le foncteur $\underline{M}\rightsquigarrow\underline{M}^{(p)}$ s'explicite
comme l'image inverse ordinaire de Modules stratifiés sur $X$ relativement à
$R$, resp.\ (dans le cas de isocristaux) de Modules stratifiés sur l'espace
rigide-analytique $X_K$, pour le « morphisme » $X_K\to X_K$ d'espaces
rigide-analytiques (relatif à $f_K$ sur le corps de base !) induit par $f_X$.
Bien que $f_X$ et par suite $f_{X_K}$ soit loin d'être unique, le foncteur
envisagé qu'il définit ne dépend pas\add{,} essentiellement\add{,} des choix
faits.

1.15. \emph{Bicristaux et biisocristaux sur une base quelconque.} Ne supposons
plus nécessairement $S$ de caractéristique déterminée $p$. Pour chaque nombre
premier $p$, soit $S_p$ la fibre $V(p.1_{\underline{O}_S})$ de $S$ sur le point
$p\underline{Z}$ de $\operatorname{Spec}(\underline{Z})$. Un bicristal de poids
$i$ sur $S$ est par définition la donnée d'un cristal de modules
$\underline{M}$ sur $S$, et pour chaque $p$ d'une structure de bi-cristal de
poids $i$ sur la restriction $\underline{M}_p$ de $\underline{M}$ à $S_p$,
définie par la donnée de $F_p,V_p$. On définit de même la notion de
bi-isocristal sur $S$, étant entendu qu'un isocristal est un objet de la
catégorie des isocristaux sur $S$ (comme de juste), définie à partir de la
catégorie des cristaux de modules en « localisant » par rapport à des entiers
$n>0$ arbitraires, i.e.\ en tensorisant les Hom sur $\underline{Z}$ par
$\underline{Q}$. -- Lorsque $S$ est de caractéristique nulle, le supplément de
structure impliqué par « bi » est évidemment vide, tandis que si $S$ est de
caractéristique $p$, on retrouve les notions de 1.10. Lorsque enfin $S$ est de
type fini sur $\underline{Z}$ et domine $\operatorname{Spec}(\underline{Z})$,
alors la notion envisagée est d'une nature essentiellement arithmétique, les
différents $F_p$ jouant le rôle d'homomorphismes de frobenius, comme de bien
entendu ; dans le cas du poids 1, en

\page{16}
\note{p.~15 de la lettre.}
particulier, on peut considérer que le cristal avec sa bi-structure
supplémentaire permet de relier entre eux les groupes formels en les diverses
caractéristiques auxquels il donne naissance \ldots

1.16. \emph{Un retour en arrière.} La définition 1.1.\ et le sorite 1.3.\ sont
un petit peu canulés. Au lieu de prendre dans 1.0.\ pour $C$ la catégorie des
\add{$R$-}flèches $S\to T$ qui \ldots, il faut prendre la catégorie des
$R$-flèches $U\to T$ qui \ldots, où $U$ est un ouvert induit non précisé de
$S$. De plus, la définition n'est guère raisonnable alors que si la catégorie
fibrée envisagée $\underline{F}$ est un « champ » pour la topologie de Zariski
sur $(\mathrm{Sch})_{/R}$, i.e.\ si on peut y recoller flèches et objets. Ceci
est nécessaire en tous cas pour pouvoir dans 1.3.\ définir la notion d'image
inverse, la définition que j'y ai donnée n'étant raisonnable que si $S,S'$ sont
affines. Dans le cas général, il faut se localiser sur $S$ et $S'$ pour se
ramener à cette situation. Autrement on bute sur des canulars idiots de nature
globale, comme le fait que sans restriction sur $S'\to S$, la construction
\struck{indiquée} \add{envisagée} de somme amalgamée fait sortir de la
catégorie des préschémas \ldots\ Il est probable qu'il y aura bien d'autres
canulars de détail dans ces notes, mais\add{,} je pense\add{,} sans
conséquence !

Il est évidemment tentant de vouloir interpréter les cristaux de modules comme
faisceaux de modules sur un certain site. C'est possible, en prenant le
« \emph{site cristallogène}\add{*} \emph{de} $S$ \emph{sur} $R$ », qui est
précisément le site dont la catégorie sous-jacente \add{$C$} est celle des
$R$-morphismes $U\to T$ ($U$ ouvert de $S$, $U\to T$ immersion fermée
surjective définie par Ideal nilpotent sur $T$), la topologie est celle de
Zariski : on prend comme familles couvrantes « de définition » de $(U\to T)$
les familles définies par des recouvrements ouverts $(T_i)$ de $T$, chaque
$T_i$ muni de la structure induite, et définissant $(U_i\to T_i)$ par
$U_i=U\cap T_i$. Un faisceau d'ensembles sur ce site s'identifie à un système
de faisceaux d'ensembles $F_{U\to T}$ sur les \add{objets} buts des objets de
$C$, avec\add{,} pour toute flèche $(U\to T)\to(U'\to T')$ de $C$\add{,} un
homomorphisme de l'image inverse de $F_{U'\to T'}$ \struck{sur dans}
\add{\uncertain{sur}} $T$, dans $F_{U\to T}$ (homomorphisme qui n'est pas
nécessairement un isomorphisme !) \add{(et qui est un isom.\ si $T\to T'$ est
une imm.\ \uncertain{ouverte})}. En particulier, associ-\note{« $R$- »,
« $C$ » et le petit « sur » au-dessus de « sur dans » biffé sont à l'encre ;
« envisagée » et « objets » sont tapés dans l'interligne et appelés à
l'encre ; l'astérisque après « cristallogène » est à l'encre ; la parenthèse
finale est ajoutée à l'encre au pied de la page, entourée et appelée après
« isomorphisme ! ) ».}

\page{17}
\note{p.~16 de la lettre.}
ant \struck{en particulier} à tout $(U\to T)$ le faisceau $\underline{O}_T$ sur
$T$, on trouve un faisceau d'anneaux $\underline{O}_C$ sur $C$. Ceci posé, les
cristaux de modules (quasi-cohérents) sur $S$ s'identifient aux Modules
quasi-cohérents (i.e.\ localement conoyau d'un homomorphisme de Modules
libres) sur $\underline{O}_C$ ; les cristaux « de \struck{type} \add{prés.}\
finie » i.e.\ les cristaux de Modules de présentation finie, correspondent
exactement aux Modules de présentation finie sur $\underline{O}_C$.\note{« en
particulier » est biffé à l'encre ; « prés. » est tapé au-dessus de « type »
biffé à la machine.}

Quand on a un $R$-morphisme $S'\to S$, je ne vois pas de morphisme naturel
correspondant entre les \emph{sites} cristallogènes correspondants à $S,S'$.
Ce n'est pas bien gênant, car introduisant les \emph{topos cristallogènes}
\add{$\mathrm{Topcr}_{S/R}$ et $\mathrm{Topcr}_{S'/R}$} \struck{définis par les
sites} en question, on définit par la méthode esquissée dans 1.3.\ un
morphisme
\[
  \mathrm{Topcr}_{S'/R}\longrightarrow\mathrm{Topcr}_{S/R}\ ,
\]
correspondant à la notion naturelle de « image inverse de faisceaux »
\add{au sens \uncertain{inverse}} (qui, j'avoue, devrait être définie avec le
plus grand soin).\note{« $\mathrm{Topcr}_{S/R}$ et $\mathrm{Topcr}_{S'/R}$ »
est tapé dans l'interligne, au-dessus de « définis par les sites », biffé à
l'encre ; « au sens inverse » (lecture douteuse du second mot) est à
l'encre, appelé avant la parenthèse.}

1.17. On est donc dans une situation où on peut faire marcher la machine
cohomologique. Etant loin de mes bases, je n'ai pas tenté sérieusement de
tirer au clair si les images directes supérieures qu'on définit ainsi, et en
particulier les groupes de cohomologie du site cristallogène, donnent ce qu'on
voudrait.\note{la fin de « cristallogène » est corrigée à l'encre sur un mot
tapé.} Le test-clef est le suivant : \emph{si $R$ est le spectre d'un corps,
et si $S$ est lisse sur $R$, et propre sur $R$ dans le cas de la
caractéristique $p>0$, on voudrait trouver, comme cohomologie à coéfficients
dans $\underline{O}_C$ lui-même, la cohomologie de De Rham de $S$
relativement à} $R$.\marginal{marche en car.~0, et \emph{pas} en car
$p>0$}\note{au crayon, dans la marge de gauche, en travers, à hauteur de ce
passage ; « pas » souligné.} Plus généralement, sans condition sur $R$, si
$f\colon S'\to S$ est propre et lisse, on voudrait trouver comme
\add{« valeur en $S$ » de} $R^if_{\mathrm{cris}\,*}(\underline{O}_{C'})$
\add{(cf 1.4)} le faisceau de cohomologie de De Rham
$\underline{R}^if_*(\underline{\Omega}^*_{S'/S})$, \struck{\ill{}} et
\add{on voudrait} que la connexion canonique à la Gauss-Manin sur ce dernier
soit celle associée à la stratification canonique provenant de la structure
cristalline (cf 1.8.).\note{« "valeur en $S$" de » et « (cf 1.4) » sont tapés
dans l'interligne et appelés à l'encre (le texte tape « $S'$ » puis le
corrige) ; « on voudrait » est à l'encre.}\marginal{(faut
\uncertain{d'introduire des} \uncertain{puiss.\ div.})}\note{au crayon, dans
la marge de gauche, à hauteur de « que la connexion canonique » ; lecture
douteuse.} L'existence d'une connexion de Gauss-Manin en cohomologie de
De-Rham (que j'ai vérifié pour n'importe quel morphisme lisse), et le tapis
de Washnitzer-Monsky, sont en tous cas des indications assez fortes pour
conduire à penser qu'il doive exister une théorie cohomologique, en termes de
cristaux, donnant de telles valeurs pour l'argument $\underline{O}_C$. Au
cas où la cohomologie du site cris-

\page{18}
\note{p.~17 de la lettre.}
tallogène ne donnerait pas les résultats qu'il faut, on pourrait essayer (au
lieu de prendre des foncteurs dérivés dans la catégorie de \emph{tous} les
Modules sur le site annelé cristallogène) de prendre des foncteurs dérivés dans
la catégorie des seuls Modules quasi-cohérents sur le site cristallogène, i.e.\
dans la catégorie des cristaux de modules.

\section*{\emph{Chapitre} 2\quad \emph{La cohomologie de De Rham est un cristal.}}

2.1. L'affirmation du titre n'est pour l'instant qu'une hypothèse ou un voeux
pieux, mais je suis convaincu qu'elle est essentiellement correcte. Comme je
l'ai dit dans 1.17., il y a pour le moment deux indications dans cette
direction :

a) Le travail de Washnitzer-Monsky. Deux grosses imperfections dans leur
tapis : 1º ils n'obtiennent que des invariants cohomologiques \emph{locaux} sur
leur variété lisse en caractéristique $p>0$, via leurs relèvements ; pour
avoir un invariant global, ils doivent se limiter à la dimension 1. On peut
penser que cela tient à leur manque de familiarité avec les machines
cohomologiques. 2º leurs invariants sont des (faisceaux de) vectoriels sur le
corps des fractions de $W$, et non sur $W$, i.e.\ ils n'obtiennent que des
invariants « modulo isogénie ». Il semble bien, en effet, que leur
démonstration d'invariance fait intervenir de façon essentielle des
dénominateurs. Il n'est pas exclu, d'après cette indication, que dans le titre
du Chapitre il faille remplacer « cristal » par « isocristal »
\add{(et « est » par « définit »)}. Ce serait bien dommage, mais n'exclurait
pas pour autant l'existence d'une bonne théorie cohomologique pour les
cristaux, qui pour un morphisme propre et lisse et le « cristal unité »
coinciderait \emph{modulo isogénie} avec ce que donne De Rham. Le test décisif
reste celui indiqué dans 1.17, savoir : donne-t-il un résultat positif
seulement en caractéristique zéro, ou en toute caractéristique ?\note{« (et
"est" par "définit") » est à l'encre, dans l'interligne.}

b) L'existence des connexions de Gauss-Manin. J'ai vérifié pour tout
morphisme lisse $f\colon X\to S$ l'existence d'une connexion canonique
\add{(absolue)} sur les $\underline{R}^if_*$ au sens de De Rham relatif, ou
plus correctement, sur le complexe\note{« (absolue) » est à l'encre ; le
$X$ de « $f\colon X\to S$ » est récrit à l'encre sur la lettre tapée, et le
« $f_*$ » complété à l'encre.}

\page{19}
\note{p.~18 de la lettre.}
$\underline{R}f_*(\underline{\Omega}^{\bullet}_{X/S})$, considéré comme objet
de la catégorie dérivée $D(\underline{O}_S)$. A vrai dire, je n'ai pas vérifié
pour cette connexion une condition de « nullité de la \add{courbure} » ;
c'est vérifié en tous cas (par voie transcendante !) si $S$ est lisse\add{,}
sur $R$ réduit de caractéristique nulle.\note{« courbure » est récrit à
l'encre sur le mot tapé.} [Il faut noter d'ailleurs qu'il n'y a pas d'espoir
de montrer que la connexion en question provient toujours d'une
stratification : c'est \emph{faux} en caractéristique $p>0$ ; la raison étant
que la cohomologie de De Rham pour une variété algébrique non complète (par
exemple affine) n'est plus du tout raisonnable, étant beaucoup trop grosse.
Pour pouvoir espérer une stratification sur
$\underline{R}f_*(\underline{\Omega}^{\bullet}_{X/S})$, sans restriction de
caractéristique, il faut donc supposer $f$ \emph{propre}.]\note{crochets
tracés à l'encre.}

c) La nécessité d'une interprétation cristallographique de la cohomologie de
De Rham est également suggérée par le problème que j'avais signalé dans mon
papier bleu sur De Rham : si $R$ est le spectre du corps des complexes, $S$
lisse sur $R$ et $X$ lisse et propre sur $S$, alors la théorie transcendante
de la cohomologie nous fournit une suite spectrale
\[
  H^p\bigl(S^{\mathrm{an}},R^qf^{\mathrm{an}}_*(\underline{C}_{S^{\mathrm{an}}})\bigr)
  \Longrightarrow H^{\bullet}(X^{\mathrm{an}},\underline{C})\ ,
\]
dont le terme initial et l'aboutissement d'autre part ont une description
purement algébrique, comme cohomologie de De Rham de $X$, et de $S$ à
coéfficients dans les $\underline{R}^qf_*(\underline{\Omega}^{\bullet}_{X/S})$
\emph{munis de leur connexion canonique} \add{intégrable} (ce qui permet de
les prendre comme coéfficients dans la cohomologie de De Rham). On aimerait
trouver une interprétation purement algébrique de cette suite spectrale toute
entière.\note{« intégrable » est tapé dans l'interligne et appelé à l'encre ;
dans la suite spectrale, $f^{\mathrm{an}}_*$ appliqué à
$\underline{C}_{S^{\mathrm{an}}}$ : sic (on attendrait
$\underline{C}_{X^{\mathrm{an}}}$).}

2.2. Je vais préciser \struck{\ill{}} l'affirmation du titre, en me plaçant
dans l'éventualité optimiste bien sûr où on n'aurait pas besoin de
s'isogéniser.\note{« bien sûr » est entouré au crayon, et un crochet au
crayon marque « dans l'éventualité » : sans doute pour déplacer les mots.}
\struck{On désigne, pour abréger, par le mot « \emph{complexe} » sur un
préschéma $S$, un complexe de Modules à cohomologie quasi-cohérente,
considéré comme objet de la catégorie dérivée qu'on sait ; Les catégories de
ces complexes, pour une base variable,}\note{passage encadré et hachuré au
crayon.} Comme les cristaux de modules sur un préschéma $S$ forment une
catégorie abélienne, on peut prendre la catégorie dérivée ; ses objets seront
appelés simplement « complexes de cristaux ». Un tel animal induit sur $S$,
et plus généralement sur tout objet-but $T$ d'un objet $(U\to T)$ du site
cristal-

\page{20}
\note{p.~19 de la lettre.}
logène \add{$C$} de $S$, un complexe de Modules ordinaire (envisagé comme
objet de la catégorie dérivée des faisceaux de Modules sur $S$, resp.\ $T$).
Un complexe de cristaux est dit pseudo-cohérent (resp.\ parfait, resp.\
\ldots) si pour tout objet $(U\to T)$ de $C$, le complexe induit sur $T$ est
pseudo-cohérent (resp.\ \ldots). Ceci posé, voilà la théorie qu'on voudrait :
A tout morphisme lisse et propre $f\colon X\to S$, serait associé un complexe
de cristaux \add{(absolu)} $\mathrm{DR}(f)$ \add{sur $S$}, appelé cohomologie
de De Rham \add{cristalline} de $f$.\note{« (absolu) » et « sur $S$ » sont
tapés dans l'interligne et appelés à l'encre, avec un petit « c » à l'encre
au-dessus de la ligne ; « cristalline » et le « $C$ » en tête de page sont à
l'encre.} Ce complexe doit être parfait, sa formation doit être compatible
avec tout changement de base sur $S$ (l'image inverse des complexes de
cristaux étant entendue, bien entendu, au sens de catégories dérivées
\ldots), et bien sûr $\mathrm{DR}(f)$ dépend fonctoriellement (de façon
cotravariante) de $X$ sur $S$. Tôt ou tard, il faudra expliciter aussi une
formule de Künneth
$\mathrm{DR}(f\times_S g)\simeq\mathrm{DR}(f)\overset{L}{\otimes}\mathrm{DR}(g)$,
et une formule de dualité, qui pour $f$ partout de dimension relative $d$
s'exprime comme un accouplement\add{,} définissant une autodualité\add{,}
\[
  \mathrm{DR}(f)\times\mathrm{DR}(f)\longrightarrow T^d[2d]\ ,
\]
où $T^d$ est le cristal de Tate de poids $2d$, et où $[2d]$ indique qu'on
translate les degrés de $2d$ (attention au facteur 2 !).\note{les crochets de
$[2d]$, le « 2 » de « poids $2d$ » et le signe $\simeq$ de la formule de
Künneth sont à l'encre.} Enfin, on veut comme de juste un isomorphisme
$\mathrm{DR}(f)(S)\simeq\underline{R}f_*(\underline{\Omega}^{\bullet}_{X/S})$,
fonctoriel en $X$, compatible avec les changements de base, avec Künneth et
la dualité (déjà connus pour la cohomologie de De Rham ordinaire).

Par prudence, je me suis abstenu de dire quoi que ce soit sur le cas $f$ non
propre, dont il faudrait parler tout au moins si on voulait faire sérieusement
le lien avec Washnitzer-Monsky.

2.3. \emph{La cohomologie de De Rham comme bi\add{iso}cristal.} A supposer
qu'on ait une théorie du type envisagé dans 2.2., \struck{\ill{}} on trouve
pour chaque entier $p$ un homomorphisme de frobenius
\[
  F_p\colon \mathrm{DR}(f)_p^{(p)}\longrightarrow\mathrm{DR}(f)_p\ ,
\]
où l'indice $p$ au complexe de cristaux $\mathrm{DR}(f)$ désigne la
restriction à $S_p=V(p.1)$ (au sens des catégories dérivées), qui d'après les
conditions de 2.2.\ n'est autre que $\mathrm{DR}(f_p)$, $f_p\colon X_p\to S_p$
étant induit par $f$. Utilisant toujours la même condi-\note{« iso » est
récrit à l'encre sur le « cr » tapé de « bicristal » ; les parenthèses autour
de « au sens des catégories dérivées » sont à l'encre. La lettre se poursuit
au lot suivant.}

\end{document}
